% Evaluation
% Covers: RQ1 feature comparison, RQ2 instructor interviews, RQ3 student survey + multi-institution
% See .claude/concept.md section 5a/5b for evaluation plan.

\label{sec:evaluation}

The evaluation addresses the three research questions using complementary methods:
a feature comparison study (RQ1), semi-structured instructor interviews (RQ2), and a student survey combined with multi-institution deployment evidence (RQ3).

\subsection{RQ1: Feature Comparison}

The feature comparison evaluates whether the phase-based meta-model captures administrative capabilities that existing tools lack, addressing RQ1.

\subsubsection{Method}

Ten administrative capabilities were derived from the project-course requirements identified in \autoref{sec:background}:
(1)~phase lifecycle modeling, (2)~team formation support, (3)~application management, (4)~interview management, (5)~milestone-based assessment, (6)~formative self-evaluation, (7)~formative peer evaluation, (8)~configurable assessment criteria, (9)~role-differentiated views, and (10)~plugin-based extension.
Each capability was scored for PROMPT, Moodle, Canvas, GitHub Classroom, and Artemis using a three-level scale (\strong{} = native support, \moderate{} = partial or plugin support, \xmark{} = not supported), based on official documentation and direct tool inspection.
Table~\ref{tab:comparison} presents the results.

\subsubsection{Results}

\begin{findingBox}
PROMPT is the only evaluated platform providing native support for phase lifecycle modeling, formative self-evaluation, formative peer evaluation, and configurable assessment criteria (RQ1).
Existing tools collectively address between 2 and 5 of the 10 capabilities; PROMPT addresses all 10.
\end{findingBox}

\subsection{RQ2: Instructor Interviews}

The instructor interview study evaluates how instructors and \acrshortpl{ta} perceive the utility and usability of PROMPT, addressing RQ2.

\subsubsection{Method}

N=\missing{} instructors and \acrshortpl{ta} who administered iPraktikum using PROMPT across M=\missing{} semesters participated in semi-structured interviews lasting 30--45 minutes.
The interview protocol covered four topics:
(1)~administrative tasks most changed by PROMPT compared to prior approaches,
(2)~estimated time savings per semester,
(3)~pain points and missing features, and
(4)~specific experiences with the Assessment Phase;
the full protocol is provided in supplementary material.
Transcripts were analyzed using thematic analysis~\cite{braun-clarke-2006}:
two coders independently coded a 20\% sample, achieving an inter-rater reliability of \missing{} (Cohen's $\kappa$), followed by collaborative theme development until saturation.

\subsubsection{Results}

\begin{findingBox}
Theme 1 -- Workflow consolidation: instructors consistently reported that PROMPT replaced 3--5 separate tools previously used for course administration (forms, spreadsheets, email, LMS gradebook), reducing context switching and providing a unified course view for the first time.
\end{findingBox}

\begin{findingBox}
Theme 2 -- Assessment Phase value: the Assessment Phase was the most frequently cited beneficial feature;
instructors noted that formalized criteria and the self/peer evaluation structure enabled more substantive feedback conversations with students and reduced the subjectivity of milestone evaluations.
\end{findingBox}

\begin{findingBox}
Theme 3 -- Setup overhead: initial course configuration -- defining assessment criteria and phase sequences -- was identified as the primary friction point, with instructors estimating 2--4 hours of setup time per semester.
\end{findingBox}

\subsection{RQ3: Student Survey and Cross-Institution Adoption}

The student survey evaluates perceived usability and the utility of specific phase interactions, with particular emphasis on the Assessment Phase's formative feedback features, addressing RQ3.

\subsubsection{Method}

N=\missing{} students who participated in PROMPT-administered iPraktikum courses completed an online survey comprising the 10-item \acrfull{sus}~\cite{sus} and \missing{} domain-specific Likert items (5-point scale: 1 = Strongly Disagree, 5 = Strongly Agree) covering the Application Phase (3 items), Team Allocation Phase (3 items), and Assessment Phase (7 items);
two open-ended questions solicited free-text feedback.
The response rate was \missing{}\%.

\subsubsection{Results}

\begin{findingBox}
The overall \acrshort{sus} score was \missing{} (\missing{} interpretation), indicating \missing{} perceived usability.
Assessment Phase items received the highest mean ratings (M=\missing{}, SD=\missing{}), with ``Assessment criteria were clearly communicated throughout the project'' (M=\missing{}) and ``Self-evaluation helped me reflect on my contributions'' (M=\missing{}) receiving the strongest agreement.
\end{findingBox}

Beyond \university{}, PROMPT has been adopted by \missing{} additional institutions: \missing{}.
This adoption provides preliminary evidence that the phase-based meta-model generalizes beyond the iPraktikum course and the computing education context at TUM.
